\begin{exercisechunk}
\item \textbf{Invariants that prevent functional completeness.}
\label[exercise]{ex:l14-boolean-invariants}
\exerciseprofile{2}{2}{2}{\exproof\quad\exinterpretation}{%
Use closure under composition to prove that two natural Boolean gate families
are not functionally complete.}
Functional completeness can fail because every available gate preserves a
property that some target Boolean function violates. This exercise develops
the monotone and affine invariants, two basic examples of that method.

For \(x,y\in\{0,1\}^r\), write \(x\le y\) when \(x_i\le y_i\) for every
\(i\). A Boolean function \(f\) is \emph{monotone} when
\(x\le y\) implies \(f(x)\le f(y)\). A Boolean function is \emph{affine over
\(\mathbb F_2\)} when it has the form
\[
 f(x)=c\oplus a_1x_1\oplus\cdots\oplus a_rx_r,
 \qquad c,a_1,\ldots,a_r\in\{0,1\},
\]
where \(\oplus\) denotes addition modulo two.
\begin{enumerate}[(a),beginpenalty=10000]
\item Prove that every circuit whose gates are monotone computes a monotone
function. Deduce that \(\{\operatorname{AND},\operatorname{OR},0,1\}\) is
not functionally complete.
\item Prove that composing affine Boolean functions produces another affine
Boolean function. Deduce that
\(\{\operatorname{XOR},\operatorname{NOT},0,1\}\) is not functionally
complete by proving that \(\operatorname{AND}\) is not affine.
\item State the general obstruction illustrated by parts~(a)--(b). Explain why
finding a composition-closed property shared by all available gates and
violated by one Boolean function is sufficient to disprove functional
completeness.
\end{enumerate}
\end{exercisechunk}
